\documentclass[10pt,a4paper]{article}
\usepackage[margin=20mm,headheight=15pt,headsep=7mm,footskip=10mm]{geometry}
\usepackage{fontspec,amsmath,amssymb,mathtools,microtype}
\setmainfont{lmroman10-regular.otf}[BoldFont=lmroman10-bold.otf,ItalicFont=lmroman10-italic.otf,BoldItalicFont=lmroman10-bolditalic.otf]
\setsansfont{lmsans10-regular.otf}[BoldFont=lmsans10-bold.otf,ItalicFont=lmsans10-oblique.otf,BoldItalicFont=lmsans10-boldoblique.otf]
\setmonofont[Scale=.86]{lmmono10-regular.otf}[ItalicFont=lmmono10-italic.otf]
\usepackage{ragged2e}\usepackage{xcolor,booktabs,longtable,array,graphicx,enumitem,fancyhdr,needspace,titlesec,hyperref,xurl}
\definecolor{AtlasNavy}{HTML}{18334C}\definecolor{AtlasTeal}{HTML}{116D73}\definecolor{AtlasGray}{HTML}{52616D}
\hypersetup{colorlinks=true,linkcolor=AtlasTeal,urlcolor=AtlasTeal,pdfcreator={Complexity Atlas},unicode=true}
\urlstyle{same}\setlength{\parindent}{0pt}\setlength{\parskip}{4.8pt}\setlength{\emergencystretch}{3em}
\renewcommand{\baselinestretch}{1.08}\setcounter{secnumdepth}{0}\setlist{nosep,leftmargin=1.5em}
\titleformat{\section}{\Large\bfseries\sffamily\color{AtlasNavy}\raggedright}{}{0em}{}
\titleformat{\subsection}{\large\bfseries\sffamily\color{AtlasTeal}\raggedright}{}{0em}{}
\titleformat{\subsubsection}{\normalsize\bfseries\color{AtlasNavy}\raggedright}{}{0em}{}
\titlespacing*{\section}{0pt}{14pt}{6pt}\titlespacing*{\subsection}{0pt}{10pt}{5pt}
\pagestyle{fancy}\fancyhf{}\fancyhead[L]{\small\sffamily\color{AtlasGray}COMPLEXITY ATLAS}
\fancyhead[R]{\small\sffamily\color{AtlasGray}ENGLISH EDITION}\fancyfoot[C]{\small\thepage}\renewcommand{\headrulewidth}{.3pt}
\newcommand{\AtlasMath}[1]{\par\begingroup\centering\setbox0=\hbox{$\displaystyle #1$}\ifdim\wd0>1.4\linewidth\typeout{ATLAS WIDE FORMULA: \the\wd0}\fi\ifdim\wd0>\linewidth\resizebox{\linewidth}{!}{\box0}\else\box0\fi\par\endgroup}
\newcommand{\AtlasInlineMath}[1]{\begingroup\setbox0=\hbox{$#1$}\ifdim\wd0>.4\linewidth\endgroup\AtlasMath{#1}\else\endgroup$#1$\fi}
\pdfstringdefDisableCommands{\def\AtlasInlineMath#1{#1}}
\DeclareMathOperator{\polylog}{polylog}\DeclareMathOperator{\dist}{dist}\DeclareMathOperator{\diam}{diam}
\DeclareMathOperator{\rank}{rank}\DeclareMathOperator{\tr}{tr}\DeclareMathOperator{\diag}{diag}
\DeclareMathOperator{\prox}{prox}\DeclareMathOperator{\KL}{KL}\DeclareMathOperator{\Var}{Var}
\providecommand{\eps}{\varepsilon}\providecommand{\R}{\mathbb R}\providecommand{\E}{\mathbb E}
\providecommand{\mathbbm}[1]{\mathbf{#1}}\providecommand{\mathds}[1]{\mathbf{#1}}
\providecommand{\llbracket}{[\![}\providecommand{\rrbracket}{]\!]}
\begin{document}
\raggedright

{\LARGE\bfseries\sffamily\color{AtlasNavy} Finite-sum convex-concave saddle-point problems: component gradients and expected global gap\par}\vspace{5pt}
\textbf{Setting ID.} \texttt{CV\_\allowbreak{}S\_\allowbreak{}FS\_\allowbreak{}SADDLE\_\allowbreak{}GAP}

Saddle-point problems

Solve \AtlasInlineMath{\min_{x\in X}\max_{y\in Y}f(x,y)}.

\textbf{Status.} Bounds missing

\textbf{TASK tags.} task family: Saddle point

\textbf{MODEL tags.} smoothness: Mean square average; structure: Finite sum

\textbf{ORACLE tags.} feedback quality: Exact; index control: Algorithm selects; interface family: Component gradient; operation: Component joint gradient; operation: Projection; public information: Hidden objects explicitly declared; public information: Public inputs explicitly declared; scalar derivative order: Scalar-function first derivative

\textbf{ALGORITHMS tags.} adaptivity: Adaptive; randomization: Private randomness allowed; restriction: Explicitly no additional algorithm restrictions; restriction: No additional span restriction explicitly imposed; restriction: No additional zero-respecting restriction explicitly imposed

\textbf{OUTPUT tags.} certificate: No additional output certificate required; object: Point

\textbf{METRIC tags.} criterion: Saddle-point gap

\textbf{GUARANTEE tags.} aggregation: Per instance; mode: Expectation; quantifier: Algorithm before instance; randomness: Algorithm private seed

\textbf{RESOURCE tags.} accounting: Joint coordinates retained; budget: Worst-instance expected cost; measure: N component; measure: N projection; setup: Initialization included

\needspace{7\baselineskip}\section{1. Setting specification}
\needspace{4\baselineskip}\subsection{Task}
Solve \AtlasInlineMath{\min_{x\in X}\max_{y\in Y}f(x,y)}.

\needspace{4\baselineskip}\subsection{Admissible model}
The set \AtlasInlineMath{K\subset\mathbb{R}^d} is known, nonempty, compact, and convex, with \AtlasInlineMath{\operatorname{diam}(K)\le D_K}. The operator \AtlasInlineMath{F=n^{-1}\sum_iF_i} is monotone; components need not be monotone.


\AtlasMath{\frac1n\sum_i\|F_i(z)-F_i(w)\|_2^2\le\beta^2\|z-w\|_2^2\quad(z,w\in K).}


Assume \AtlasInlineMath{K=X\times Y}; the average \AtlasInlineMath{f=n^{-1}\sum_i f_i} is convex-concave, and \AtlasInlineMath{F_i=(\nabla_xf_i,-\nabla_yf_i)} holds. Individual \AtlasInlineMath{f_i} need not be convex-concave.

\needspace{4\baselineskip}\subsection{Oracle interface}
Public inputs include block dimensions \AtlasInlineMath{d_x,d_y} and \AtlasInlineMath{D_K}, \AtlasInlineMath{\beta}, \AtlasInlineMath{\varepsilon}, and \AtlasInlineMath{n}, together with the stated initial points and feasible-set information. Analytic representations of hidden functions, operators, and matrices are not supplied.


\AtlasMath{(i,x,y)\longmapsto(\nabla_xf_i(x,y),\nabla_yf_i(x,y))}


The algorithm selects an index and feasible point. A component may be queried at multiple points, with each evaluation charged. Values are exact; randomness comes from the algorithm, not the oracle.

Exact Euclidean projection onto the known set is allowed. Its calls are recorded separately and excluded from the primary query count.

\needspace{4\baselineskip}\subsection{Output}
Return \AtlasInlineMath{\widehat{z}\in K}, equivalently \AtlasInlineMath{(\widehat{x},\widehat{y})\in X\times Y}.

The output may be any measurable postprocessing of the history, including an unqueried point, average, or convex combination. It must satisfy the stated domain and feasibility requirements; no additional function value or computable certificate is required.

\needspace{4\baselineskip}\subsection{Loss or accuracy criterion}

\AtlasMath{\max_{v\in Y}f(\widehat{x},v)-\min_{u\in X}f(u,\widehat{y})}


Compactness makes this gap finite; expectation is applied to the complete gap.

\needspace{4\baselineskip}\subsection{Resource accounting}
One full vector response from a selected component costs one \AtlasInlineMath{N_{\rm component}} call. A full average costs \AtlasInlineMath{n} calls. Projection counts are retained separately.

Use expected work on the worst instance. The algorithm stops almost surely, with call count \AtlasInlineMath{\tau} satisfying \AtlasInlineMath{\sup_I\mathbb{E}[\tau]\le N}. This does not require every execution path to finish within \AtlasInlineMath{N} calls.

Count all information calls for one instance and one requested accuracy, including initialization. Exact-real arithmetic, storage, and solving known models from public inputs and returned data are uncharged. This query-complexity quantity does not certify arithmetic or bit complexity.

\needspace{4\baselineskip}\subsection{Algorithm views and success guarantees}
Arbitrary measurable sequential adaptive randomized algorithms are allowed. The private seed is independent of the hidden instance. Queries, stopping, and output may use the full history, without additional span, zero-respecting, communication, or memory restrictions.

Choose one algorithm, then require the guarantee for every admissible instance and stochastic-feedback mechanism.


\AtlasMath{\mathbb{E}\!\left[\max_{v\in Y}f(\widehat{x},v)-\min_{u\in X}f(u,\widehat{y})\right]\le\varepsilon.}


Expectation is outside the complete error expression; internal maxima or suprema are not moved outside it.

\textbf{Budget convention.} worst\_instance\_expected\_calls

\needspace{7\baselineskip}\section{2. Notation and parameters}
\AtlasInlineMath{D_K}: A Euclidean-diameter bound on the entire feasible product domain. \textbf{Domain:} Positive real number.

\AtlasInlineMath{\operatorname{diam}(K)=\sup_{u,v\in K}\|u-v\|_2\le D_K}

\AtlasInlineMath{\beta}: Mean-square average Lipschitz upper bound for component joint gradients \textbf{Domain:} Positive real number.

\AtlasInlineMath{\frac1n\sum_{i=1}^n\|\nabla f_i(z)-\nabla f_i(w)\|_2^2\le\beta^2\|z-w\|_2^2}

\AtlasInlineMath{d}: The joint Euclidean dimension is \AtlasInlineMath{d=d_x+d_y}. Within its stated window, the lower bound may choose finite block dimensions. \textbf{Domain:} Integer at least \AtlasInlineMath{2}.

\AtlasInlineMath{\varepsilon}: Target tolerance for the complete weak-VI gap or global saddle gap, in the units of that gap. \textbf{Domain:} Positive real number.

\AtlasInlineMath{n}: Fixed number of finite-sum components; the full function or operator is their 1/n average \textbf{Domain:} Positive integer.

\AtlasInlineMath{\|\cdot\|_2,\langle\cdot,\cdot\rangle}: Euclidean norm and inner product; the product-space norm is the square root of the sum of the squared block norms. \textbf{Domain:} The specified finite-dimensional real vector space.

\AtlasInlineMath{\|v\|_2=\sqrt{\langle v,v\rangle}}

\AtlasInlineMath{d_x,d_y}: Positive integer primal and dual block dimensions; joint dimension \AtlasInlineMath{d=d_x+d_y} \textbf{Domain:} Positive integer.

\AtlasInlineMath{f}: Unknown real-valued payoff, with primal and dual input blocks. \textbf{Domain:} \AtlasInlineMath{f:(x,y)\mapsto f(x,y)\in\mathbb R}.

\AtlasInlineMath{x,y,z}: Primal variable, dual variable, and joint point \textbf{Domain:} \AtlasInlineMath{x\in\mathbb R^{d_x},\ y\in\mathbb R^{d_y},\ z=(x,y)\in\mathbb R^{d_x+d_y}}.

\AtlasInlineMath{z=(x,y)}

\AtlasInlineMath{\nabla f,\nabla_x f,\nabla_y f}: Complete joint gradient and its two blocks \textbf{Domain:} \AtlasInlineMath{\nabla_x f\in\mathbb R^{d_x},\ \nabla_y f\in\mathbb R^{d_y}}.

\AtlasInlineMath{\nabla f=(\nabla_x f,\nabla_y f)}

\AtlasInlineMath{f_i,i}: The fixed ith real-valued component and the algorithm-selected index. \textbf{Domain:} \AtlasInlineMath{i\in\{1,\ldots,n\},\quad f_i:\text{the declared product domain}\to\mathbb R}.

\AtlasInlineMath{f=\frac1n\sum_{i=1}^n f_i}

\AtlasInlineMath{K}: A known nonempty compact convex feasible set. \textbf{Domain:} A subset of Euclidean space.

\AtlasInlineMath{K=X\times Y\subset\mathbb R^{d_x+d_y}}

\AtlasInlineMath{X,Y}: Both primal and dual domains are known, nonempty, compact, and convex. Their product has diameter at most \AtlasInlineMath{D_K}. \textbf{Domain:} \AtlasInlineMath{X\subset\mathbb R^{d_x},\ Y\subset\mathbb R^{d_y}}.

\AtlasInlineMath{\widehat x,\widehat y}: The point or variable blocks ultimately returned by the algorithm, corresponding to the hatted variables in the error formula. \textbf{Domain:} Feasible variable pairs in \AtlasInlineMath{\mathbb R^{d_x}\times\mathbb R^{d_y}}.

\AtlasInlineMath{\widehat z}: Joint output consisting of the two variable blocks \textbf{Domain:} \AtlasInlineMath{\mathbb R^{d_x+d_y}}.

\AtlasInlineMath{\widehat z=(\widehat x,\widehat y)}

\AtlasInlineMath{\mathbb E,\Pr}: Expectation and probability over the explicitly allowed algorithm and sample randomness, with the hidden instance fixed. \textbf{Domain:} Operations on the corresponding probability space.

\AtlasInlineMath{N_{\rm component}}: Number of full joint-gradient queries to a specified component \textbf{Domain:} Nonnegative integer count; expected work is its expectation.

\AtlasInlineMath{N_{\rm projection}}: Projection count, uncharged in the primary query-cost projection \textbf{Domain:} Nonnegative integer count; expected work is its expectation.

\AtlasInlineMath{\tau}: Cumulative primary information calls at termination of one run; almost surely finite, with random stopping allowed. \textbf{Domain:} A nonnegative integer-valued random variable.

\AtlasInlineMath{\sup_I\mathbb E[\tau]\le N}

\AtlasInlineMath{N}: A worst-instance upper bound on expected total information calls at termination. \textbf{Domain:} A nonnegative real.

\AtlasInlineMath{N^\star}: The infimum principal query budget among algorithms satisfying every stated instancewise output guarantee. \textbf{Domain:} Nonnegative extended real value; by convention, positive infinity if no successful algorithm or no uniform finite budget exists.

\AtlasInlineMath{\mathcal A,I}: An algorithm selected from public inputs and an admissible hidden instance fixed before interaction \textbf{Domain:} The specified algorithm and model classes.

\AtlasInlineMath{\omega}: A complete realization of algorithmic and sample randomness. \textbf{Domain:} Corresponding sample space.

\AtlasInlineMath{O,\Omega,\Theta,\widetilde O}: Asymptotic upper and lower bounds, constant-factor matching, and upper bounds hiding the stated logarithmic factors, within the corresponding regions. \textbf{Domain:} Nonnegative complexity function; subscript p permits additional dependence only on fixed p.

\AtlasInlineMath{F}: Saddle operator induced by the payoff function \textbf{Domain:} \AtlasInlineMath{X\times Y\to\mathbb R^{d_x+d_y}}.

\AtlasInlineMath{F(x,y)=(\nabla_xf(x,y),-\nabla_yf(x,y))}

\AtlasInlineMath{F_i}: Saddle-point operator associated with the \AtlasInlineMath{i}th component gradient \textbf{Domain:} \AtlasInlineMath{X\times Y\to\mathbb R^{d_x+d_y}}.

\AtlasInlineMath{F_i=(\nabla_xf_i,-\nabla_yf_i)}

\needspace{12\baselineskip}\section{3. Complexity bracket}
\needspace{4\baselineskip}\subsection{Compared quantities}
\needspace{4\baselineskip}\subsubsection{Q1 - source scoped metric resource statement}
\textbf{Metric contract.} output: domain: zhat in K; postprocessing: arbitrary measurable transcript-based output unless algorithm restriction expressly states otherwise; metric: max\_\{y in Y\} f(xhat,y)-min\_\{x in X\}f(x,yhat); guarantee: One algorithm family chosen from public inputs satisfies for every legal instance and oracle law: E[max\_\{y in Y\} f(xhat,y)-min\_\{x in X\}f(x,yhat)]<=epsilon.

\textbf{Resource contract.} budget: sup\_instance E[N\_oracle] (expected stopping-work budget); primary: N\_component; unit: one declared oracle call; auxiliary: projection and known prox evaluations separately retained

This is the native source metric/resource guarantee. Joint precision, block size, runtime and oracle components stay separate; do not infer a scalar optimal query complexity.

\needspace{4\baselineskip}\subsection{R1 - The bound’s native applicability window.}
\textbf{Region.} Assume \AtlasInlineMath{n\ge2} and \AtlasInlineMath{\varepsilon>0}. Both total component work and the global saddle-gap guarantee are in expectation.

\textbf{Quantity.} Q1

\textbf{Resource.} Joint resource budget

\textbf{Status.} Lower bound missing

\textbf{Lower bound.} No admitted bound for this comparison

\textbf{Upper bound.} B1

\needspace{7\baselineskip}\section{4. Bounds and supporting evidence}
The following conditions apply to every bound in this profile, in addition to each bound's stated region and algorithm scope.

\textbf{Conditions.} Use the original algorithm’s uniform certificate inequality together with the convex-concave saddle-point ergodic inequality, then take the supremum over comparison points to control the global gap. A small weak VI gap at a point alone does not imply a small saddle gap there.

\needspace{4\baselineskip}\subsection{B1 - See the cited theorem}
\textbf{Admission.} Included in the bracket

\textbf{Direction.} upper

\textbf{Witness or lower-bound scope.} Randomized witness

\textbf{Algorithm-scope qualification.} The variance-reduced extragradient/Mirror-Prox algorithm of Alacaoglu-Malitsky genuinely samples components and refreshes reference points randomly.

\AtlasMath{O\!\left(n+\frac{\sqrt n\,\beta D_K^2}{\varepsilon}\right)}
\textbf{Sources.} \href{https://proceedings.mlr.press/v178/alacaoglu22a/alacaoglu22a.pdf}{[1]}

\textbf{Constants.} Use the source theorem’s universal constants except where derivative order or structural budgets are explicitly fixed. Unspecified sufficient thresholds remain symbolic.

\textbf{Source location.} \href{https://proceedings.mlr.press/v178/alacaoglu22a/alacaoglu22a.pdf}{[1]} Assumption1; Algorithm1; Theorem4; Remark5; Corollary6

Expected total component calls O(n+sqrt(n)L D\_K\textasciicircum{}2/epsilon) for expected Minty weak gap.

Source writes sum F, ours average: source components F\_i/n, uniform estimator returns our F\_i. Same index evaluated at two points, two charged calls. No individual component monotonicity. Full pass n. Native loopless count is expected work. Saddle-gap transfer uses source Eq12-14 and convex-concave inequality, not equivalence of output Minty gap and full saddle gap.

\needspace{7\baselineskip}\section{References}
[1] Ahmet Alacaoglu; Yura Malitsky. Stochastic Variance Reduction for Variational Inequality Methods. 2022. COLT2022 PMLR178.

\url{https://proceedings.mlr.press/v178/alacaoglu22a/alacaoglu22a.pdf}

\end{document}
